\subsection{RANK OF AN ELEMENT OF A TENSOR PRODUCT}

Let E be a right vector space and F a left vector space over the same field K; to every element \(u \in E \otimes_{K} F\) there corresponds canonically under (29) a homomorphism \(u_{1} \in \operatorname{Hom}_{\mathbb{K}}(\mathbf{E}^{*}, \mathbf{F})\) ; if \(u = \sum_{i} x_{i} \otimes y_{i}\) with \(x_{i} \in E, y_{i} \in F\) , the element \(u_{1}\) is the linear mapping

\[
x ^ {*} \mapsto \sum_ {i} \langle x ^ {*}, x _ {i} \rangle y _ {i}.\tag{33}
\]

On the other hand, \(\mathbf{E} \otimes_{\mathbb{K}} \mathbf{F}\) is canonically identified with \(\mathbf{F} \otimes_{\mathbb{K}^0} \mathbf{E}\) , where \(\mathbf{E}\) is considered as a left vector space and \(\mathbf{F}\) as a right vector space over the opposite field \(\mathbf{K}^0\) ; thus there corresponds canonically to \(u\) a homomorphism \(u_2 \in \operatorname{Hom}_{\mathbb{K}}(\mathbf{F}^*, \mathbf{E})\) given by

\[
y ^ {*} \mapsto \sum_ {i} x _ {i} \langle y _ {i}, y ^ {*} \rangle ;\tag{34}
\]

\(u_{1}\) (resp. \(u_{2}\) ), considered as a mapping of \(E^{*}\) into \(F^{**}\) (resp. of \(F^{*}\) into \(E^{**}\) ) is just the transpose of \(u_{2}\) (resp. \(u_{1}\) ). The ranks of \(u_{1}\) and \(u_{2}\) are thus equal to the same finite number r, the common dimension of the subspaces \(u_{1}(E^{*})\) of F and \(u_{2}(F^{*})\) of E, each of which is canonically isomorphic to the dual of the other (no. 5); we shall say that r (denoted \(\operatorname{rg}(u)\) ) is the rank of the element u of \(E \otimes_{K} F\) and that \(u_{1}(E^{*})\) and \(u_{2}(F^{*})\) are the subspaces (of F and E respectively) associated with u.

PROPOSITION 17. Let \(u\) be an element of \(\mathbf{E} \otimes_{\mathbb{K}} \mathbf{F}\) and \(\mathbf{M} \subset \mathbf{E}\) and \(\mathbf{N} \subset \mathbf{F}\) its associated subspaces. For every expression \(u = \sum_{i=1}^{s} x_i \otimes y_i\) of \(u\) , where \(x_i \in \mathbf{E}\) and \(y_i \in \mathbf{F}\) for \(1 \leqslant i \leqslant s\) , the subspace \(\mathbf{M}\) (resp. \(\mathbf{N}\) ) is contained in the subspace of \(\mathbf{E}\) (resp. \(\mathbf{F}\) ) generated by the \(x_i\) (resp. the \(y_i\) ). Moreover, the following properties are equivalent:

\begin{enumerate}
\def\labelenumi{(\alph{enumi})}
\item
  The integer \(s\) is equal to the rank of \(u\) .
\item
  The family \((x_{i})_{1\leqslant i\leqslant s}\) is a basis of M.
\item
  The family \((y_{i})_{1\leqslant i\leqslant s}\) is a basis of N.
\item
  The families \((x_{i})_{1\leqslant i\leqslant s}\) and \((y_{i})_{1\leqslant i\leqslant s}\) are both free.
\end{enumerate}

By (33) (resp. (34)) each element of \(N = u_{1}(E^{*})\) (resp. \(M = u_{2}(F^{*})\) ) is a linear combination of the \(y_{i}\) (resp. the \(x_{i}\) ); whence the first assertion. If s = r, the subspace generated by the \(x_{i}\) (resp. \(y_{i}\) ) with dimension \(\leqslant \dim M\) (resp. \(\leqslant \dim N\) ) and containing M (resp. N) is identical with it and hence (a) implies (b) and (c) and a fortiori (d). Conversely, each of conditions (b) and (c) implies (a) by definition of \(\operatorname{rg}(u)\) . Finally if (d) holds, there exists a family \((x_{i}^{*})_{1\leqslant i\leqslant s}\) of elements of \(E^{*}\) such that \(\langle x_{i}, x_{j}^{*}\rangle = \delta_{ij}\) (no. 5, Corollary 1 to Theorem 7) and hence it follows from (33) that \((y_{i})\) is a basis of N, which completes the proof.

COROLLARY 1. The rank of \(u\) is the smallest integer \(s\) such that there exists an expression \(u = \sum_{i=1}^{s} x_i \otimes y_i\) , where \(x_i \in \mathbf{E}\) and \(y_i \in \mathbf{F}\) for \(1 \leqslant i \leqslant s\) .

This follows immediately from Proposition 17 and no. 2, Proposition 1.

COROLLARY 2. Let K be a commutative field, E, F two vector spaces over K and L a commutative extension field of K. Let u be an element of E ⊗K F, M and N its associated subspaces and u' the canonical image of u in (E ⊗K F)(L) (canonically identified with E(L) ⊗L F(L), cf.~§ 5, no. 1, Proposition 3); then rg(u') = rg(u) and the subspaces associated with u' are canonically identified with M(L) and N(L).

If \(u = \sum_{i=1}^{r} x_i \otimes y_i\) , where the families \((x_i)\) and \((y_i)\) are free, then \(u' = \sum_{i=1}^{r} (1 \otimes x_i) \otimes (1 \otimes y_i)\) and the families \((1 \otimes x_i)\) and \((1 \otimes y_i)\) are free in \(E_{(L)}\) and \(F_{(L)}\) respectively (§ 5, no. 1, Proposition 4).

